% !TEX root = ../auv-thesis.tex
% Authoritative thesis source; edit this TeX file directly.
% Legacy Markdown material is reference-only and must not overwrite this file.

\chapter{绪论}

\section{研究背景与重大意义}\label{sec:ch01-1-1}

\subsection{海上风电的战略地位}\label{sec:ch01-1-1-1}

要理解本文的研究对象为何重要，首先需要把视野放回到国家能源结构调整的整体背景中。在过去十余年里，海上风电正在经历一场从``近海示范''到``深远海大规模布局''的结构性迁移：单机容量从早期的 3 MW 级跨越到 16 MW 级以上，场站规模从单期数十台扩展到百万千瓦级集群，离岸距离也从 10 km 量级推向 50--100 km 甚至更远。这种空间和容量的同时扩张，使得海上风电不再是沿海地方性能源补充，而是国家能源外送骨架的关键组成部分。

在这种发展格局下，海底电缆扮演着不可替代的角色。它是连接发电端（风机)、升压站（海上变电平台）和陆上电网三者的物理通道，也是海上风电从``发出''到``消纳''的唯一传输路径：无论上游风机多么稳定、下游电网调度多么先进，只要海缆中断，整条能量传输链即被打断。因此，海缆的可靠性并不仅仅是设备层面的问题，而是直接影响区域电力供应安全和可再生能源消纳能力的系统性问题。

在我国``双碳''目标推进的背景下，深远海风电已进入国家新能源布局。国家九部门印发的《``十四五''可再生能源发展规划》提出优化近海海上风电布局、开展深远海海上风电规划，并推动近海规模化开发和深远海示范；随后发布的新能源消纳政策进一步把海缆廊道、登陆点与海上输电网络的集约化布局纳入海上风电基地开发的系统约束\cite{ndrc2022renewableplan}。规模侧的官方发布亦说明这种约束正在快速显性化：``十四五''期间我国海上风电新增装机超过 3500 万千瓦\cite{scio2025energyachievement}；从全球视角看，2024 年海上风电累计装机已达 83 GW，中国占据约一半市场份额\cite{gwec2025offshore}。因此，本文不将海缆仅理解为风机之间的附属连接，而将其视为需要持续获得路由、埋深与运行状态证据的关键能源基础设施；``海缆资产快速增长、巡检需求同步放大''正是本文研究 AUV 巡检的直接问题背景。

\subsection{海缆：风电场的``生命线''}\label{sec:ch01-1-1-2}

如果把整个海上风电场视作一个系统，那么风机相当于``心脏''，海上升压站相当于``枢纽''，而海底电缆则相当于把心脏的输出送到外部的``主动脉''。这个比喻虽然朴素，却比``重要附件''更能说明海缆在场站中的地位：风机停一台，损失的是单台容量；海缆断一条，损失的可能是整片场站的全部出力。这种``单点失效、全场受影响''的属性，决定了海缆运维必须按``关键基础设施''的标准对待，而不是按一般电气设备的常规巡修对待。

进一步看，海缆故障的成因相当多样：船舶锚损是最常见的人为破坏因素，海床冲刷会逐步暴露原本掩埋的电缆使其更易受外力作用，潮汐和水流变化会引起悬跨段的疲劳损伤，台风、地震和滑坡等自然灾害则可能在很短时间内造成大段电缆位移甚至断裂。面向海上可再生能源海缆的损伤与寿命预测研究也表明，机械载荷、海床条件和环境作用必须联合纳入风险评估\cite{dinmohammadi2019predicting}。每一种成因都对应一种不同的早期征兆，掩埋深度变浅、路由发生偏移、屏蔽层局部异常、磁场或温度场出现异动，这些征兆若能被巡检设备及早捕捉到，就有可能避免一次完整的停电事故。

海缆故障的经济后果不能用单一的``故障率''概括，因为修复船、备缆与海况窗口共同决定停运时间。DNV 引述 Lloyd Warwick 的保险损失资料显示，海缆故障关联约 83\% 的海上风电财务损失与保险索赔；其给出的阵列缆和送出缆平均修复停运时间分别约为 40 d 和 60 d，修复加发电损失的合并损失区间分别为 120 万-\/-1200 万美元和 1000 万-\/-3000 万美元\cite{dnv2024cablefailure}。这些数据是保险与损失口径而非普适故障率，却足以说明早期发现异常的价值：诊断、定位与维修若被推迟，损失会沿着海况窗口和船舶资源约束放大。

正是在这种现实需求下，本节强调的不是``海缆需要被巡检''，而是``海缆巡检必须在故障真正发生之前就开始周期性进行''。这也与监管侧对运行期的要求一致：上海海事局规定业主单位应定期监测海缆路由及附近海域水下地形\cite{shanghai2022windtraffic}。因此，后续章节的状态估计、决策架构和控制策略，最终服务的都是在故障发展为停运之前发现埋深、路由或运行状态征兆这一目标。

\subsection{运维现状与挑战}\label{sec:ch01-1-1-3}

在深入讨论 AUV 之前，有必要先把``目前是怎么巡检的''这件事说清楚。当前海缆巡检在工程实践中主要采用人工潜水、母船拖曳式探测、ROV 和有人驾驶水面/水下平台等方式。国内方法选择研究和近期海洋机器人综述均将声学、电磁、ROV 与 AUV 视为按任务条件组合的技术谱系，而非彼此完全替代的单一路线\cite{cen2017submarine,zhang2025submarine}。海上风电海缆运维的近期综述表明，现有体系已覆盖状态监测、检测与全生命周期风险评估，但仍需把设备状态、海床环境和运维资源统筹为可执行的运维策略\cite{ning2026cableom}。这些方式各自有合适的使用场景，但放在``深远海大规模周期性巡检''的需求下，都存在不同程度的局限。

人工潜水的优势是对小范围、近岸场地的精细观察具有不可替代性，但在水深和距离限制下不可能成为主力手段；母船拖曳式探测可在较大范围内提供连续的声学数据，但拖鱼姿态受船速和缆索张力影响，姿态不稳定时探测精度会显著下降；ROV 是当前覆盖能力最强的工具，可操作性强、可实时干预，对故障点确认和精细维修具有不可替代性，但其作业必须依托母船、缆控连接和专业操作团队，单日作业成本动辄数十万元，在长期、周期性巡检任务中很难做到经济可持续。

把这些局限叠加到``深远海''这一新约束上，问题进一步放大：母船在深远海作业受海况和补给约束，单次出海窗口有限；缆控 ROV 受脐带缆、母船定位和操作团队制约，不易经济地覆盖长距离路由；人员安全风险也随离岸距离与任务持续时间上升。ORE Catapult 与 RenewableUK 的行业报告将水面、水下和空中机器人平台列为未来海上风电运维的重要能力，明确其目标是降低运维成本并改善作业安全\cite{orecatapult2025ras}。因此，深远海的关键矛盾并非简单以一种新装备替代人工，而是如何让不同平台在覆盖、确认、维修与安全保障之间形成可交接的作业链。

AUV 正是在这种压力下被引入的。它的核心价值不在于``能做 ROV 做的所有事''，而在于：可以在预设任务约束下自主航行、近底扫描、持续记录多源传感数据、在通信不稳的条件下保持任务连续性，并明显降低对母船的实时依赖。这意味着同一个母船每次出海可以放出多台 AUV 并行作业，每台 AUV 在水下持续 8--24 小时进行高密度扫描，使单位面积巡检成本和单位时长安全风险相应下降。

但是，AUV 也并非万能。本文在定位时刻意避免``AUV 全面替代 ROV''的过度叙述，而采用更工程化的分工定位：AUV 适合大范围、周期性、前置筛查式巡检，通过覆盖整段路由发现可疑征兆、给出待复核的目标清单；ROV 仍然是后续精细维修和故障确认的核心工具，根据 AUV 的筛查结果进入指定区域执行高精度作业。本文研究的所有内容都建立在这种``AUV 前置筛查 + ROV 后置确认''的分工假设之上。

\section{海底管缆探测技术研究现状}\label{sec:ch01-1-2}

\subsection{声学探测技术}\label{sec:ch01-1-2-1}

声学是水下探测最早成熟、也是目前最主流的感知手段。围绕海缆巡检任务，常用的声学设备主要包括两类：侧扫声呐（Side-Scan Sonar，SSS）和浅地层剖面仪（Sub-Bottom Profiler，SBP）。它们的物理原理都是发射声脉冲、接收海底反射并根据回波时延和强度构建图像，但二者关注的``层次''不同，侧扫声呐关注海底表面的纹理和裸露目标的几何形状，浅地层剖面仪则关注海床浅层的内部结构。

把这两种声学手段套到海缆巡检任务上，其用途也很自然：当电缆部分裸露在海底时，侧扫声呐能给出比较直观的纹理和走向线索；当电缆被掩埋在浅层沉积物之下时，浅地层剖面仪可以通过低频声波的穿透性给出垂向反射剖面，从而推断电缆位置和大致埋深。这是声学方法在海缆巡检中的``理想工况''。

然而，把``理想工况''换成真实海底环境，问题就开始出现。在浑浊水域，悬浮颗粒会增加散射、降低有效分辨率；在底质变化剧烈的区域（例如砂泥过渡带、岩石碎屑带），声学反射特征本身就足够复杂，使电缆特征被淹没在背景中；当电缆掩埋深度超过声波穿透能力时，浅地层剖面仪的回波将无法可靠区分电缆与沉积层；声呐图像的解译目前仍高度依赖人工经验，自动识别算法在复杂底质下的鲁棒性远未成熟。

正因为存在这些局限，声学探测虽然是必不可少的``看得到几何形状''的工具，但单独使用难以应对掩埋、屏蔽、底质突变等复杂工况。这些局限本身，就是本文后续引入``声磁多源融合''的最直接动机：当声学信号不稳定时，系统需要另一个互补的物理通道来维持对电缆位置和状态的估计能力。

侧扫声呐用于 AUV 海缆/管线自主跟踪已有系统性研究\cite{bagnitsky2011sidescan,feng2023automatic}，多平台声学建图也被用于海缆路由勘测\cite{jung2022multimodal}；针对掩埋目标，浅地层剖面仪凭借低频穿透在管线与海缆探测中发挥作用\cite{geomatching2022sbp,zheng2021zeroshot}，三维声学成像更被用于完全掩埋风电海缆的探测\cite{ha2022buried}。

\subsection{电磁感应探测技术}\label{sec:ch01-1-2-2}

与声学方法基于``反射几何''不同，电磁感应方法的物理基础是``电流产生磁场''。海缆中流动的工频电流会在周围空间形成一个周期性磁场，借助通量门磁强计（Fluxgate）或磁感应式传感器，可以测量到这个磁场在传感器位置的三维分量，并通过反演方法估计电流走向、相对距离和大致埋深。

这种方法在物理上具有几个声学方法不具备的优势。第一，磁场不依赖光学或声学透明度，因此对浑浊度不敏感，浊水中仍可正常探测；第二，磁场可在一定厚度的非磁性介质中穿透，对掩埋电缆比浅地层剖面仪更具直接探测能力；第三，磁场探测能直接给出``是否有运行电流''这一二值证据，对识别``电缆存在但已停运''和``电缆运行中''非常关键。这些特性使电磁感应方法成为声学的天然互补手段。

然而，三相交流海缆引入了一个明显的难点。三相电流的相位差为 120°，使得三相电流在某些方向上的磁场会相互抵消，从而显著降低有效漏磁信号的强度；铠装层（金属护套）会通过涡流屏蔽进一步衰减外部可测磁场；导线的螺旋绞合则使漏磁特征在空间上呈现周期性而非单调变化。这三种效应叠加，意味着即使在电缆正下方，磁场强度也可能比单芯直流电缆弱出一两个数量级。这就要求完整测量链具备 0.05 nT 量级的系统磁异常分辨能力，并配套更精细的标定流程。

电力行业对此提出了明确要求。例如《DL/T 1278-2025 海底电力电缆运行规程》对磁探测灵敏度、埋深反演误差和路由精度都有具体条款，本项目以 0.05 nT 级系统磁异常分辨能力和 0.2 m 级埋深反演误差为设计目标。这些数值不是任意选定的，而是来源于行业规程对探测精度的硬性约束；前者需由传感器、采集链和窄带处理联合验证，而不是单个器件的手册指标。

总结来说，电磁感应方法的优势在于``看得到电流''，劣势在于``信号弱、易受干扰、三相抵消''，而其与声学方法的互补特性，正是后续多源融合方案得以成立的物理基础。

在磁探测传感与算法层面，已有工作涵盖高灵敏度隧道磁阻（tunnel magnetoresistance, TMR）传感器的结构优化\cite{bi2025tmr}、AUV 磁导引海缆跟踪\cite{xiang2016subsea,zhang2023adaptive}、基于标量磁强计的姿态无关路由跟踪\cite{li2024attitude}，以及水下矢量磁数据处理\cite{seidel2023uxo,oehler2024processing}与地表磁场探测输电电缆\cite{sun2014underground}等。近期研究又分别从弱磁交流建模\cite{su2025precise}和近岸工频磁特征航磁定位\cite{sun2026magnetic}两个方向推进了运行海缆的路由识别。本项目的 0.05 nT 灵敏度与 0.2 m 埋深误差指标则以《DL/T 1278---2025 海底电力电缆运行规程》\cite{dlt1278_2025}和《海底管道及海底电缆检测技术规范》\cite{tcsgpc_cable}为直接依据，三相海缆漏磁的定量分析可进一步参照三芯海缆损耗解析模型\cite{hatlo2015accurate}。

\subsection{多源感知融合技术}\label{sec:ch01-1-2-3}

如果声学方法擅长``看几何''、电磁方法擅长``看电流''，那么把二者结合起来似乎是顺理成章的选择。但工程实践中，多源融合并不是把两个传感器的数据相加那么简单，它面临一系列具体挑战：声学数据通常以图像帧或扫描行为单位，时间戳颗粒度大；磁数据则是连续时间序列，需要高频采样；二者的空间参考点不同（声呐在 AUV 下方挂点、磁传感器在长杆末端），需要严格的杆臂改正；二者的观测质量随环境波动剧烈，例如声呐在浊水中失效时，磁数据可能正常，反之亦然。

更核心的问题在于：现有大多数研究在做声磁融合时，默认两个传感器都``基本可靠''，因此融合算法的核心是设计加权或最优估计。但真实海况下的常态恰恰相反，某段时间内声学失效、某段时间内磁场被外部干扰、再某段时间内多普勒测速仪（Doppler velocity log, DVL）丢底导致整体运动估计退化。这种``传感器轮流不靠谱''的情形，使得``是否融合''和``如何加权''成为比``加权多少''更关键的问题。

本文据此提出一个不同的视角：把``不确定性感知''作为多源融合的前提，而不是假设传感器始终可靠。具体来说，融合系统不仅要计算每个观测对状态估计的修正量，还要在线评估每个观测的可信度，并把这种可信度作为权重的一部分；当某通道可信度低于阈值时，系统应自动降低对该通道的依赖，甚至在极端情况下完全不使用该通道。这种思想贯穿后续状态估计（第 3 章）和控制设计（第 4 章），并在不确定性感知模型预测控制（uncertainty-aware model predictive control, UA-MPC）的代价函数中得到直接体现。

水下组合导航的技术路线也从基于模型的 EKF/UKF 对比\cite{allotta2015comparison}逐步扩展到数据驱动的深度序列导航\cite{zhang2020navnet}。前者具有状态和协方差可解释、便于纳入传感器物理模型的优势，后者则能学习复杂时序误差，但依赖训练分布和数据覆盖。本文选择 ES-EKF 作为在线融合骨架，不是排斥学习方法，而是因为论文当前更需要可追溯地表达观测退化、协方差传播和安全降级条件。

\subsection{典型探测装备与自主化沿革}\label{sec:ch01-1-2-4}

海缆巡检装备的演进需要区分三个层次：第一层是直接产生路由和埋深观测的探测载荷，第二层是搭载载荷的拖体、ROV、AUV 或水面平台，第三层才是包含任务规划、目标跟踪、异常处置和数据闭环的自主作业系统。把三者混为一谈容易得出``某型探测仪就是自主巡检机器人''的错误结论。现有综述也表明，工程系统通常是在具体任务约束下组合探测原理、运载平台与作业流程\cite{zhang2025submarine,cen2017submarine}。

在专用探测载荷方面，Teledyne TSS 的产品谱系体现了从单一电缆跟踪向多模式管缆探测的发展。HydroPACT 350 面向带有源音频信号的海缆，通过两组三轴线圈给出相对位置、埋深、载体偏航角和前视信息；其产品规格给出的低噪声条件探测范围为垂向 10 m、以线圈阵列为中心的水平总宽度 20 m，测量精度取 5 cm 与斜距 5\% 中的较大值\cite{teledyne2021hydropact350}。HydroPACT 440 采用脉冲感应探测导电目标，660E 进一步面向小型观察级 ROV 减小尺寸和载荷重量，Dualtrack 则把 350 与 440 组合以覆盖有源音频和脉冲感应两类工况\cite{teledyne2026hydropactline}。Innovatum SMARTRAK 则在单套设备中集成被动磁、有源交流和有源直流三种模式，可安装于 ROV、AUV、拖体或水面船\cite{innovatum_smartrak}。这些系统说明，商用探测仪的核心竞争力首先是稳定输出可测量的横向偏移和埋深，而其自主程度取决于搭载平台是否把观测接入在线制导闭环。

国内装备已形成专用探测仪和自主载体两条并行路线。海鹰 T180 采用有源电磁原理，提供 25 Hz、50 Hz 和 133 Hz 工作频率以及路由、埋深和故障测量模式，其公开指标包括 200 m 最大使用水深、路由误差为水深的 \(\pm 5\%\)、埋深误差为埋深的 \(\pm 10\%\)\cite{haiying_t180}。在工程应用侧，HydroPACT 350 已在琼州海峡福康线 500 kV 带电海缆上完成路由与埋深试验，说明运行海缆可利用自身工频特征开展不停电测量，但现场仍需处理风机等外部工频干扰\cite{teledyne2017hainan350}。在自主平台侧，华中科技大学深圳研究院公开的飞行式长艏翼样机面向海缆配置双三轴正交电磁探测设备，标称深度等级为 300 m\cite{hust2023cableauv}；精灵 P200 则以模块化 AUV 集成导航、多波束和成像声呐，公开指标为 300 m 潜深、3 kn 下续航 12 h，并明确包含海底电缆巡检应用\cite{fsea_p200}。这些资料反映出国产装备已从``可搭载探测载荷''推进到``面向海缆场景的专用样机''，但公开的长期海试里程、复杂底质检出率和自主故障处置证据仍有限。

国外自主化路线则从通用 AUV 测绘逐步走向资产级跟踪与岸基长航程作业。REMUS 系列早期工作奠定了小型 AUV 模块化任务载荷与自主航行的工程基础\cite{purcell2000remus}；HUGIN 系列进一步把合成孔径声呐、多波束和自主管缆跟踪集成到长航程平台，形成从近海到 6000 m 级深水的多型号体系\cite{kongsberg2025hugin}。面向海上风电海缆，Anduril DIVE-LD 项目已开展岸基布放的全自主示范，任务包括 Point Judith-\/-Block Island 海缆路由的侧扫声呐与测深广域扫描，以及对冲刷防护和异常目标的近距离复核\cite{nowrdc2024diveld}。这条路线的实质不是简单取消脐带缆，而是把覆盖路径规划\cite{choset2001coverage}、目标重捕获、质量控制和异常复核一并下沉到载体侧。

图~\ref{fig:ch01-prior-equipment} 从实物角度对照了上述装备演进的三个代表：专用有源电磁探测载荷通过多组感应线圈直接测量海缆漏磁（图~\ref{fig:ch01-hydropact350}、图~\ref{fig:ch01-t180}），自主载体则把此类载荷与导航、声呐一并集成到可长航程巡检的 AUV 平台上（图~\ref{fig:ch01-p200}）。这也说明本文关注的重点不在复刻单一载荷，而在把观测接入可自主决策的巡检闭环。

\begin{figure}[htbp]
\centering
\subcaptionbox{Teledyne HydroPACT 350 双三轴感应线圈组\cite{teledyne2021hydropact350}\label{fig:ch01-hydropact350}}{%
  \includegraphics[width=0.31\linewidth,height=0.2\textheight,keepaspectratio,alt={HydroPACT 350 两组三轴感应线圈与黑色电子舱}]{literature/prior_work_equipment/hydropact350_coils.png}}\hfill
\subcaptionbox{海鹰 T180 海缆探测系统探测棒\cite{haiying_t180}\label{fig:ch01-t180}}{%
  \includegraphics[width=0.31\linewidth,height=0.2\textheight,keepaspectratio,alt={海鹰 T180 黄色潜水与水面探测棒线圈阵列及连接缆}]{literature/prior_work_equipment/haiying_t180.jpg}}\hfill
\subcaptionbox{精灵 P200 模块化自主水下机器人\cite{fsea_p200}\label{fig:ch01-p200}}{%
  \includegraphics[width=0.31\linewidth,height=0.2\textheight,keepaspectratio,alt={精灵 P200 橙色鱼雷形自主水下机器人置于水面浮排上}]{literature/prior_work_equipment/fsea_p200.png}}
\caption{典型海缆探测载荷与自主载体实物示例：从专用有源电磁探测线圈组到集成式自主水下机器人}\label{fig:ch01-prior-equipment}
\end{figure}
\FloatBarrier

\begingroup\scriptsize

\begin{longtable}[]{@{}
  >{\raggedright\arraybackslash}p{(\linewidth - 8\tabcolsep) * \real{0.2000}}
  >{\raggedright\arraybackslash}p{(\linewidth - 8\tabcolsep) * \real{0.2000}}
  >{\raggedright\arraybackslash}p{(\linewidth - 8\tabcolsep) * \real{0.2000}}
  >{\raggedright\arraybackslash}p{(\linewidth - 8\tabcolsep) * \real{0.2000}}
  >{\raggedright\arraybackslash}p{(\linewidth - 8\tabcolsep) * \real{0.2000}}@{}}
\caption{典型海缆探测载荷与自主巡检平台的能力分层}\label{tab:ch01-cable-inspection-equipment}\tabularnewline
\toprule\noalign{}
层级 & 代表装备 & 主要观测/能力 & 典型搭载或作业方式 & 对本文的启示 \\
\midrule\noalign{}
\endfirsthead
\toprule\noalign{}
层级 & 代表装备 & 主要观测/能力 & 典型搭载或作业方式 & 对本文的启示 \\
\midrule\noalign{}
\endhead
\bottomrule\noalign{}
\endlastfoot
专用有源电磁载荷 & HydroPACT 350、海鹰 T180 & 运行或注入音频下的路由、横向偏移和埋深 & ROV、拖体、水面平台 & 工频窄带提取与横向偏移反演应形成稳定接口 \\
脉冲感应/多模式载荷 & HydroPACT 440、660E、Dualtrack；SMARTRAK & 导电目标、被动磁、有源 AC/DC 多模式探测 & ROV、AUV、拖体 & 不同物理通道需要质量标志和模式切换机制 \\
通用自主调查平台 & REMUS、HUGIN、精灵 P200 & 声呐测绘、组合导航、长航程覆盖 & 船载或岸基布放 AUV & 自主能力不能只用``无缆''定义，还需任务闭环证据 \\
海缆专用自主样机/示范 & 长艏翼海缆巡检 AUV、DIVE-LD & 电磁跟踪或声学广扫加异常复核 & 近底自主跟踪、岸到岸任务 & 本文应聚焦观测可信度到决策控制的闭环，而非复刻单一商用品 \\
\end{longtable}

\endgroup

综上，典型商用品已经较好解决``怎样稳定测到电缆''的问题，先进 AUV 则在解决``怎样以更少母船依赖持续测量''的问题；两者之间仍缺少的是在声学、磁学和导航观测轮流退化时，如何维持可解释的在线估计并安全调整任务。本文的声磁协同 ES-EKF、行为树和 UA-MPC 正是针对这一系统层缺口，而不是宣称在单项探测距离或商用成熟度上替代上述设备。

\section{AUV 自主控制与决策系统现状}\label{sec:ch01-1-3}

\subsection{轨迹跟踪控制}\label{sec:ch01-1-3-1}

AUV 的运动控制是整个系统能否完成巡检任务的物理基础。从控制理论的角度看，AUV 是一个典型的强非线性、欠驱动、受流场扰动的水下刚体；从工程实现的角度看，控制器还要兼顾算力约束、传感器延迟和实时性要求。围绕这种被控对象，主流方法包括比例-\/-积分-\/-微分控制（proportional-\/-integral-\/-derivative control, PID）、线性二次型调节（linear quadratic regulator, LQR）、视线（line-of-sight, LOS）制导和模型预测控制（model predictive control, MPC）。其中，比例 LOS 的稳定性分析\cite{fossen2014usges}与 LOS-MPC 路径跟踪\cite{oh2010path}构成了从几何制导向预测优化过渡的代表性基础。

PID 是工程上使用最广泛的方法，因其实现简单、调参直观且在线性时不变近似下足够稳定。在 AUV 的深度环、航向环等单维度控制任务上，调优良好的 PID 通常能给出令人满意的结果，本文采用的 Python 船舶仿真器（Python Vehicle Simulator, PVS）内置自动驾驶仪即基于这种思想。但 PID 的局限也很清晰：它逐通道独立调节，难以显式处理执行器饱和、状态约束、复杂路径预瞄等多变量耦合问题，对不确定性的适应能力依赖经验式参数整定。

LQR 在一定程度上缓解了 PID 的多通道耦合问题，通过最小化二次型代价获得最优反馈律。但 LQR 的最优性建立在线性化模型和无约束假设之上，对 AUV 这种强非线性、欠驱动系统来说，要么需要大量分段线性化、要么必须借助增益调度，工程复杂度增加而最优性优势减弱。

MPC 引入了``预测时域 + 在线优化''的思路，可以在每个控制周期显式处理状态约束、控制约束和优化目标。AUV 领域已从路径规划与跟踪的统一滚动优化\cite{shen2017integrated}，发展到模型辨识结合预测控制的试验验证\cite{zhou2023parameter}以及规划-\/-跟踪一体化设计\cite{deng2024motion}。理论上，MPC 比 PID 和 LQR 更适合复杂路径跟踪、速度规划、多目标权衡等任务，并且天然支持把感知不确定性映射到代价函数中；随机 MPC 综述则系统说明了不确定性传播、概率约束与在线计算之间的基本权衡\cite{mesbah2016stochastic}。其代价是对动力学模型精度和在线求解算力有更高要求，部署到嵌入式平台需要在求解精度、求解时间和稳定性之间小心权衡。

那么，这三种方法在本文的任务中分别处于什么位置？本文的思路是按任务特性区分使用：在近底地形跟随这种``目标变化平缓、强调安全裕度''的任务中，调优后的 PID/PVS 作为当前最可靠的方案；在复杂平面路径制导任务中，MPC 的预瞄和速度规划开始展现优势；在不确定性场景中，UA-MPC 的作用不是保证横向精度全面提升，而是把感知置信度转化为控制风格和跟踪许可，使系统在高压力代理场景中保持可解性、降低控制变化率，并避免把低质量观测误用为连续跟踪证据。本文据此采取``按任务特性选择最合适的控制层级''的策略，而不是追求``一种控制器覆盖所有任务''；相应的实验证据将在第 5 章给出，分层策略的设计细节则在第 4 章展开。

\subsection{自主决策架构}\label{sec:ch01-1-3-2}

如果说控制器决定``AUV 怎么走得稳''，那么决策架构决定``AUV 决定走哪儿、什么时候停下来、出问题怎么办''。在自主水下机器人领域，主流的决策架构有两种：有限状态机（finite state machine, FSM）和行为树（behavior tree, BT）。

FSM 是较早期的决策范式，其核心思想是把任务划分为一组离散状态，并在状态之间定义明确的迁移条件。FSM 在状态数较少时简单清晰，是早期自主车辆和早期 AUV 的常见选择。但当任务复杂度上升、状态数显著增加时，FSM 的几个问题就会暴露出来：状态间的迁移条件呈``状态数的平方''增长，维护成本急剧上升；新增一种异常处理常常需要在所有相关状态中添加迁移边，难以局部修改；不同优先级行为之间的抢占关系需要显式建模，逻辑容易出现遗漏。

行为树通过``控制节点 + 条件节点 + 行动节点''的层级组织，提供了一种更模块化的决策表达方式。控制节点（Sequence、Selector、Parallel）以组合子的方式描述子节点的执行语义，条件节点检查系统状态，行动节点执行具体行为。这种结构天然支持任务可组合（一个新任务可作为子树插入）、异常可抢占（高优先级 failsafe 节点放在 Selector 顶端）、扩展可局部化（新增一个分支不影响其他分支）。对于像 AUV 巡检这样需要``正常任务执行 + 异常抢占自愈''并存的任务，行为树明显比 FSM 更合适。

水下机器人中的行为树应用已从多智能体覆盖任务的安全约束组合\cite{ozkahraman2020combining}，扩展到欠驱动 AUV 的非线性 MPC 与行为树联合控制\cite{bhat2023controlling}和协同水下作业的动态任务规划\cite{choi2025dynamic}。更近期的机器人故障处置研究还尝试把反应式规划、行为树和视觉语言模型结合起来\cite{ahmad2025unified}。本文不引入视觉语言模型，而是吸收这些工作的共同原则：安全行为应具有最高抢占优先级，故障处理应以可验证的局部子树接入主任务。

本文据此选用行为树作为自主决策的核心框架，把搜索、锁定、跟踪、应急上浮、失联保护等行为组织在同一个决策树中。具体的节点设计、分支组织和与控制层的接口将在第 4 章展开。

\section{本文研究内容与核心创新点}\label{sec:ch01-1-4}

\subsection{研究内容}\label{sec:ch01-1-4-1}

本文面向海上风电海底电缆 AUV 自主巡检中``弱磁目标感知不确定、近底环境复杂、传感器模态间歇失效''三类核心困难，提出声磁协同状态估计、感知不确定性驱动控制与行为树安全决策相结合的分层虚实验证框架，并通过多场景、多噪声、多失效模式、多控制器消融及真实 PC104 零执行器实验，验证其在设定条件下的有效性与适用边界。围绕这一目标，本文的研究内容涵盖以下四个层次，从底层硬件到上层任务策略，构成一条完整的工程链条：

在算法性能验证之外，本文还以真实 PC104/VxWorks、生产 ROS2 仲裁器和行为树完成零执行器同步安全链验证；因此下文所称``纯仿真''只限定算法性能矩阵，不再用于概括整篇论文的证据层级。

\begin{enumerate}
\def\labelenumi{\arabic{enumi}.}
\item
  \textbf{面向 AUV 海缆巡检的软硬件系统集成}：构建包含执行层（AMD Geode LX 800 实时控制板）、感知决策层（Jetson/第二代机器人操作系统，Robot Operating System 2, ROS2）、上位机和仿真后端的完整架构，支持从算法开发、仿真验证到实物部署的渐进式工程路径。这部分研究关注架构选型、协议设计、安全闭环和虚实切换接口。
\item
  \textbf{声磁协同状态估计与不确定性量化}：在 DVL 丢包、声呐杂波、磁信号畸变等多源不确定性共存的环境中，维持状态估计的连续性，并把感知不确定性以可量化的方式传递给后续控制层。这部分研究关注误差状态扩展卡尔曼滤波（error-state extended Kalman filter, ES-EKF）的算法设计、观测协方差在线调整和置信度衰减机制。
\item
  \textbf{近底巡检任务的决策与控制策略}：以行为树完成任务调度，PID/PVS 作为当前最可靠的近底地形跟随主控制器，制导级 MPC（包括 UA-MPC 变体）作为复杂路径预瞄和不确定性感知扩展。这部分研究关注分层控制架构、安全仲裁和在不确定性下的控制权重调整。
\item
  \textbf{分层虚实结合验证体系与适用边界界定}：通过 PVS 高保真仿真、HoloOcean 视觉验证、故障注入、地形跟随基准与真实 PC104 零执行器同步实验形成可复现的验证闭环，并对当前结果的统计充分性、场景真实性和硬件实物证据三类边界给出明确说明。这部分研究关注实验设计、指标定义，以及在分层验证阶梯上把系统级计量与真机验收组织为可独立核查的递进层次。
\end{enumerate}

把这四部分组合起来，可以覆盖``硬件 → 估计 → 控制 → 验证''的完整链路，每一层的设计选择都为下一层提供约束和接口。

图~\ref{fig:ch01-technical-route} 将上述工作组织为一条由任务约束驱动、经系统集成与算法设计、最终落到验证边界的技术路线。其中，状态估计输出的置信度是连接感知与决策控制的关键共享信号。

\begin{figure}[htbp]
\centering
\includegraphics[width=0.95\linewidth,height=\textheight,keepaspectratio,alt={任务约束驱动双脑系统、声磁估计和分层决策控制，并由多层验证界定迁移边界}]{architecture/auv_thesis_technical_route.pdf}
\caption{面向海底电缆自主巡检的技术路线：任务约束驱动双脑系统、声磁估计和分层决策控制，并由多层验证界定迁移边界}\label{fig:ch01-technical-route}
\end{figure}
\FloatBarrier

\subsection{核心创新点}\label{sec:ch01-1-4-2}

基于本文已完成的工作，可以梳理出以下三类相对扎实的创新点，它们分别对应主线中的``声磁协同状态估计''、``感知不确定性驱动控制''与``行为树安全决策''三个环节：

\begin{enumerate}
\def\labelenumi{\arabic{enumi}.}
\item
  \textbf{面向海缆弱磁目标的声磁协同状态估计仿真框架}：针对三相交流海缆螺旋漏磁的相互抵消、钢丝铠装层的弱漏磁衰减等使目标磁信号微弱的物理特征，利用三轴磁场模量的旋转不变性与之字形横切形成的半高全宽几何关系反演电缆垂直距离，降低对未知电流幅值、铠装衰减和传感器姿态的依赖；并以 ES-EKF 为骨架融合惯性、DVL、声学与磁场异步观测，对电缆相对位置、AUV 状态与目标置信度进行联合估计，使观测协方差随信噪比、遮挡、声呐失效和磁异常在线自适应调整。
\item
  \textbf{设计并实现了感知不确定性驱动的模型预测控制（UA-MPC）}：不确定性不作为事后评价，而是通过 ES-EKF 协方差到置信度的映射函数进入控制环——把观测质量和状态估计不确定性转化为可被控制层直接消费的标量信号，并在 UA-MPC 中映射为跟踪权重的非线性放大和控制代价的 sigmoid 平滑调整，使观测可靠时优先提高跟踪精度、观测退化时主动降低控制激进程度。可解性、最大迭代与高机动三组消融分别从可解性、迭代上限和横向机动精度三方面验证了该机制，表明控制代价降权是主要贡献，而单纯增加最大迭代次数不能构成实时控制修复。
\item
  \textbf{面向自主巡检任务的行为树安全决策与分层验证体系}：以行为树作为任务级安全调度机制，在搜索、捕获、沿线跟踪、失效恢复与安全悬停/返航等行为间按高优先级抢占切换，并以分层任务状态机与应急逻辑保证安全兜底；通过故障注入、传感器失效、目标丢失等仿真场景与真实 PC104 故障字同步实验，以反应延迟、抖动、恢复时间和任务完成率等指标验证其在故障下真实起作用，而非仅为流程图式的工程实现。作为支撑，执行层与感知决策层的双脑解耦架构（轻量二进制 UDP、任务下发通道与上位机授权）使仿真验证与实物部署共享上层算法、仅替换接口层即可切换运行环境，为上述安全逻辑提供工程载体。
\end{enumerate}

上述创新点均限定于仿真与数字孪生口径：磁反演一致性由三轴复 I/Q 联合拟合支持，UA-MPC 的 PVS 消融支持可解性与控制平滑归因，确定性运动学对照支持横向精度归因，行为树安全逻辑由故障注入对照支持。其中行为树的 Bit5/Bit13/Bit14 状态传递另有真实 PC104--ROS2 同步零执行器证据，把安全状态链的实证等级提升到实物侧。这些结论尚不替代真实检测噪声、真机实时性与硬件计量条件下的实物验收，也不把仿真控制性能外推为执行机构或整艇水域性能。

\section{本文组织架构}\label{sec:ch01-1-5}

本文按照``背景 → 系统 → 算法 → 控制 → 实验 → 总结''的逻辑展开，每一章为下一章提供必要的概念和约束：

\begin{itemize}
\tightlist
\item
  \textbf{第 1 章}阐述研究背景、技术现状和本文工作，建立``为什么要做这件事''的问题意识；
\item
  \textbf{第 2 章}介绍任务需求分析、双脑硬件架构、声磁物理建模和标定框架，回答``系统由什么组成、为何这样组成''；
\item
  \textbf{第 3 章}阐述声磁协同状态估计算法、不确定性量化模型和仿真环境下的验证方案，回答``系统如何感知自己和环境''；
\item
  \textbf{第 4 章}介绍行为树任务调度、轨迹生成策略和不确定性感知控制设计，回答``系统如何决定动作和执行控制''；
\item
  \textbf{第 5 章}汇总实验平台、评价指标、已有结果、讨论与 Sim-to-Real 迁移性分析，回答``我们已经验证了什么、还差什么''；
\item
  \textbf{第 6 章}总结全文的声磁协同估计、不确定性感知控制、近底安全约束与执行链证据，并讨论系统级计量、真机闭环与海试迁移等剩余验证方向。
\end{itemize}
